% Copyright (c) 2026 Eugene Joseph M. Ragasa.
% SPDX-License-Identifier: Apache-2.0
% Author: Eugene Joseph M. Ragasa
% Formal proof completed: 2026-11-02
\documentclass[11pt]{article}
\usepackage[T1]{fontenc}\usepackage{amsmath,amssymb}\usepackage[margin=1in]{geometry}\usepackage[hidelinks]{hyperref}
\title{Proof Dossier: Normalized FCC Complete-Shell Gauge-Support Radius}\author{Eugene Joseph M. Ragasa}\date{Formal proof completed: 2026-11-02}
\begin{document}\maketitle
\section{Identity and status}
ID \texttt{BG-FOURIER-SUPPORT-THM-016}; status \texttt{proved}. Lean declaration:
\par\noindent{\tiny\path{ProjectKoios.BandGauge.Fourier.Support.}}\\
{\tiny\path{fccShellRadius_shell_neg_add_add_subset}}.

\section{Exact radius}
For integer reduced coordinates $m=(x,y,z)$, the existing normalized FCC shell index is
\[
 q(m)=x^2+y^2+z^2+xy+xz+yz.
\]
Define the additive Euclidean realization
\[
 v(m)=\frac{1}{\sqrt 2}(x+y+z,x,y,z)
\]
and the nonnegative radius $\rho(m)=\lVert v(m)\rVert_2$. Direct expansion proves
$\rho(m)^2=q(m)$ exactly. Because $v(-m)=-v(m)$ and $v(m+n)=v(m)+v(n)$, the Euclidean norm gives
\[
 \rho(-m)=\rho(m),\qquad \rho(m+n)\leq\rho(m)+\rho(n).
\]
These prerequisites are proved separately:
\begin{itemize}
\item \texttt{BG-FOURIER-SUPPORT-THM-004} gives the exact squared-index identity;
\item \texttt{BG-FOURIER-SUPPORT-LEM-006} gives inversion invariance; and
\item \texttt{BG-FOURIER-SUPPORT-THM-008} gives the triangle inequality.
\end{itemize}

\section{Gauge--field--gauge bound}
Suppose two gauge modes $u_1,u_2$ obey $\rho(u_i)\leq R_U$ and a field mode $a$ obeys
$\rho(a)\leq R_A$. Inversion invariance and two applications of the triangle inequality give
\[
 \rho(-u_1+a+u_2)
 \leq \rho(u_1)+\rho(a)+\rho(u_2)
 \leq R_A+2R_U.
\]
The elementwise theorem \texttt{BG-FOURIER-SUPPORT-THM-012} specializes a carrier-independent
helper to the normalized FCC radius. Defining the complete shell
$S_\rho(R)=\{m\mid \rho(m)\leq R\}$ then gives the dossier theorem
\[
 (-S_\rho(R_U)+S_\rho(R_A))+S_\rho(R_U)
 \subseteq S_\rho(R_A+2R_U).
\]
Membership in the two pointwise set sums supplies $u_1,a,u_2$; the elementwise estimate proves
membership in the shell on the right.

\section{Boundary}
This is a support inclusion on the infinite integer lattice, not a support equality. It does not identify a
finite Fourier character with a centered integer representative. A finite periodic application must
separately prove quotient and seam compatibility, including modular wraparound. Cancellation may
make actual coefficient support smaller. No runtime mesh, FFT convention, coefficient array,
physical length, hopping range, material applicability, or scientific adequacy is certified.

\section{Trusted base and verification}
Checked with Lean \texttt{4.35.0} and Mathlib revision
\texttt{56a005855ea17e0b0958f73ec801ae140b0a1679}. Kernel axioms are
\texttt{propext}, \texttt{Classical.choice}, and \texttt{Quot.sound}; no repository-owned axiom or
admission is used. AI-assisted and kernel-checked; no independent expert review recorded.
\end{document}
